\clearpage
\begin{examplebox}
\textbf{例1：内心题里的竖直角平分线}

三角形 $PAB$ 的内心为 $I$，已知 $P(2,1)$，$I$ 在 $x=2$ 上，$PA,PB$ 均不竖直。求两边的斜率关系。
\begin{center}
\begin{tikzpicture}[x=1cm,y=1cm]
\draw[p055 frame] (0,0) rectangle (13.8,7.8);
\node[anchor=west,font=\small\bfseries] at (.4,7.3) {例1：边射线互为镜像，不要求两条边等长};
\begin{scope}[shift={(3.3,1.45)},scale=.78]
  \coordinate (P) at (2,6); \coordinate (A) at (0,3); \coordinate (B) at (6,0);
  \coordinate (Ap) at (4,3); \coordinate (I) at (2,{5+(5-9*sqrt(13))/(5+3*sqrt(13))});
  \draw[p055 mirror] (2,-.3)--(2,6.5);
  \node[teal!70!black,left] at (2,.5) {$x=2$};
  \draw[gray] (A)--(B);
  \draw[p055 original] (P)--(A);
  \draw[p055 image] (P)--(B);
  \draw[p055 helper] (A)--(Ap);
  \fill (P) circle (2pt) node[above left] {$P(2,1)$};
  \fill[CStatement] (A) circle (2pt) node[left] {$A(0,-2)$};
  \fill[orange!85!black] (B) circle (2pt) node[right] {$B(6,-5)$};
  \fill[orange!85!black] (Ap) circle (2pt) node[right] {$A'(4,-2)$};
  \fill (I) circle (2pt) node[left] {$I$};
  \draw[gray] (P)--(I);
  \node[anchor=west,align=left] at (7.1,4.3) {$A'$ 在射线 $PB$ 上\\[5pt]$PA=PA'$\\[5pt]$PB=2PA$};
  \node[anchor=west,align=left] at (7.1,1.5) {$I$ 在内角平分线上\\[5pt]$k_{PA}=\frac32,\ k_{PB}=-\frac32$};
\end{scope}
\node at (6.9,.45) {虚线连接真正互为镜像的点 $A,A'$；$B$ 只需落在像射线上。};
\end{tikzpicture}

\end{center}
\textbf{解：} 内心在内角平分线上，故 $PI$ 平分 $\angle APB$；又 $PI$ 为竖直线 $x=2$，所以两边射线互为反射，得到
\[
\boxed{k_{PA}=-k_{PB}}.
\]
\textbf{回验：} 图取 $A(0,-2),B(6,-5)$，斜率依次为 $3/2,-3/2$。$A$ 的镜像是 $A'(4,-2)$，它在射线 $PB$ 上；$PA=\sqrt{13},\ PB=2\sqrt{13}$，两边确实无需等长。此三角形的内心横坐标为 $2$（用三边长加权坐标也可核验）。

\medskip
\textbf{例2：水平镜面}

$P(-1,4),A(0,2)$，将射线 $PA$ 关于 $y=4$ 翻折，求像斜率。

\textbf{解与回验：} 原增量 $(1,-2)$，水平翻折后为 $(1,2)$，故 $k'=2$。对称点 $A'=(0,6)$，直接计算 $(6-4)/(0+1)=2$，一致。
\end{examplebox}

\clearpage
\begin{examplebox}
\textbf{例3：$45^\circ$ 镜面，既求方向又求位置}

$P(0,1),Q(1,3)$，将直线 $PQ$ 关于 $y=x+1$ 反射，求像直线。
\begin{center}
\begin{tikzpicture}[x=1cm,y=1cm]
\draw[p055 frame] (0,0) rectangle (13.8,6.7);
\node[anchor=west,font=\small\bfseries] at (.4,6.2) {例3：先换方向，再用镜面上的不动点定位};
\begin{scope}[shift={(2.1,1.3)},scale=1.05]
  \draw[p055 axis] (-.5,0)--(4.1,0) node[right] {$x$};
  \draw[p055 axis] (0,-.4)--(0,3.6) node[above] {$y$};
  \draw[p055 mirror] (-.6,.4)--(2.5,3.5) node[right] {$y=x+1$};
  \draw[p055 original] (-.5,0)--(1.25,3.5);
  \draw[p055 image] (-.6,.7)--(4,3);
  \draw[p055 helper] (1,3)--(2,2);
  \fill (0,1) circle (1.8pt);
  \node[anchor=east] at (-.15,1.55) {$P(0,1)$};
  \fill[CStatement] (1,3) circle (1.8pt);
  \node[CStatement,above] at (1,3.5) {$Q(1,3)$};
  \fill[orange!85!black] (2,2) circle (1.8pt) node[below right] {$Q'(2,2)$};
  \node[anchor=west,align=left] at (6.1,2.6) {$k=2\ \longmapsto\ k'=\frac12$\\[7pt]$P$ 在镜面上，保持不动\\[7pt]像直线：$y=\frac12x+1$};
\end{scope}
\node at (6.9,.45) {截距影响镜面位置；它不进入斜率变换 $k'=1/k$。};
\end{tikzpicture}

\end{center}
\textbf{解：} 原斜率 $k=2$，像斜率 $k'=1/2$；$P$ 在镜面上，保持不动，所以像直线为 $\boxed{y=\frac12x+1}$。独立求得 $Q'=(2,2)$，代入像方程成立。

\medskip
\textbf{例4：一般镜面，通式算完再反查}

把 $r:y=1$ 关于 $l:y=2x+1$ 反射，求像直线。

\textbf{解：} $m=2,k=0$，故
\[
k'=\frac{2m+(m^2-1)k}{1-m^2+2mk}=\frac4{-3}=-\frac43.
\]
原线与镜面交于不动点 $P(0,1)$，所以像直线为 $\boxed{y=-\frac43x+1}$。

\textbf{独立回验：} $Q(5,1)$ 到镜面的垂足是 $H(1,3)$，故 $Q'=2H-Q=(-3,5)$；$(5-1)/(-3-0)=-4/3$，一致。
\end{examplebox}

\clearpage
\begin{examplebox}
\textbf{例5：反射后竖直，不能硬算有限斜率}

将 $r:y=\frac34x+1$ 关于 $y=2x+1$ 反射，求像直线。
\begin{center}
\begin{tikzpicture}[x=1cm,y=1cm]
\draw[p055 frame] (0,0) rectangle (13.8,8.3);
\node[anchor=west,font=\small\bfseries] at (.4,7.8) {例5：新横增量为零，像直线仍然存在};
\begin{scope}[shift={(2.6,1.25)},scale=.86]
  \draw[p055 axis] (-.7,0)--(5,0) node[right] {$x$};
  \draw[p055 mirror] (-.4,.2)--(2.7,6.4) node[right] {$y=2x+1$};
  \draw[p055 original] (-.4,.7)--(4.7,4.525);
  \draw[p055 image] (0,.1)--(0,6.5);
  \draw[p055 helper] (4,4)--(0,6);
  \fill (0,1) circle (2pt);
  \node[anchor=east] at (-.2,1.6) {$P(0,1)$};
  \fill[CStatement] (4,4) circle (2pt) node[below right=4pt] {$Q(4,4)$};
  \fill[orange!85!black] (0,6) circle (2pt) node[left] {$Q'(0,6)$};
  \node[anchor=west,align=left] at (6.2,4.7) {原增量 $(4,3)$\\[6pt]像增量 $(0,5)$\\[6pt]$\Delta x'=0$，$\Delta y'\ne0$};
  \node[anchor=west,orange!85!black] at (6.2,2.1) {像直线：$x=0$};
\end{scope}
\node at (6.9,.45) {有限斜率不存在 $\ne$ 像直线不存在。};
\end{tikzpicture}

\end{center}
\textbf{解：} 通式分母为 $1-4+4\times\frac34=0$，所以像直线竖直。又 $P(0,1)$ 在镜面上，反射后不动，因此像直线为 $\boxed{x=0}$。

\textbf{回验：} $Q(4,4)$ 的垂足为 $H(2,5)$，对称点 $Q'=(0,6)$ 与 $P$ 同在 $x=0$ 上；图中原增量长为 $5$，像增量长也为 $5$。

\medskip
\textbf{例6：斜率反号，仍不能确定指定镜面}

$y=x$ 与 $y=-x+2$ 斜率为 $1,-1$，是否关于 $x=0$ 对称？

\textbf{解：} 不是。真正的像是 $y=-x$；题设两线交于 $(1,1)$，该点不在镜面上。下一节的位置对照图把这两条下降直线同时画出，直接显示差别。
\end{examplebox}
